\section{Reusable Continuation in a Family of Capital Decisions}
\label{sec:r16menu}

The learned Bellman object can be reused when current economic decisions
change but their continuation does not. We examine that property within
the capital economy. The finite-period experiment changes the current
weight on consumption, the current adjustment cost, or a temporary
consumption constraint. Every subsequent primitive, innovation law,
terminal preference, and reference policy is common across these decisions.
The stored numerical coefficients define the finite economy exactly;
the experiment does not identify a quarterly Bellman problem with a
continuous-time diffusion without a further approximation argument.

Write the first postdecision mean as
$x_a(y)=y+h\{c\bm1+f(y)-a\}$. For a current task $\theta$, let
\[
 Q_\theta(y,a)=r_\theta(y,a)+S\{x_a(y)\},\qquad
 S(x)=\E F(x;\xi),
\]
where $F$ contains the current innovation and all subsequent reference
rewards, excluding the current reward. The finite query catalog, task
catalog, action boxes, and complete distribution over training streams
are fixed before confirmation.

\begin{proposition}[The common continuation object]
\label{prop:r16menucache}
If two current tasks differ only in their current reward or current
feasible set, their postdecision continuation $S$ is identical. A cached
continuation observation at the same $x$ and the same innovation array
is therefore valid for either task. Cached innovations may also be
reused when $x$ changes, but the corresponding state recursion and payoff
must then be reevaluated. A cached payoff or derivative at a different
$x$ is not an exact observation at the new argument.
\end{proposition}

\begin{proof}
Conditional on $x$, every future transition kernel, reward, terminal
payoff, and reference action is the same. Iterating these common kernels
gives the same law and payoff functional, and hence the same expectation.
The last assertion follows because the transition recursion depends on
its starting state even when its innovations are fixed.
\end{proof}

This distinction determines the comparison. NBO fits one scalar
continuation to common rollout labels and queries its derivative for
different current choices. The vector-costate fit, cached Raw actor,
task-conditioned direct actor, and cached sample-average optimizer are
all allowed their own legitimate reuse. In particular, the Raw competitor
does not regenerate an innovation array merely because an optimizer
revisits a state. The experiment compares the economic decisions and the
work required by these complete procedures. It does not infer that every
form of continuation reuse requires a neural scalar critic.

\begin{proposition}[A fixed-dictionary critic's bias and variance]
\label{prop:r16menuprojection}
Fix postdecision arguments and positive weights before drawing a new
costate-label array. Let $q$ be the vector of true continuation gradients,
after multiplication by the square roots of these weights, and let
$z=q+\varepsilon$, with $\E\varepsilon=0$ and finite covariance
$\Sigma$. Fix differentiable scalar features independently of this
label array, and let $P$ be the Euclidean orthogonal projection onto
their stacked, weighted gradient columns. The least-squares scalar
critic has fitted gradients $Pz$, and
\[
 \E\|Pz-q\|^2-\E\|z-q\|^2
   =\|(I-P)q\|^2-\operatorname{tr}\{(I-P)\Sigma\}.
\]
Thus its risk reduction equals the label variance removed outside the
gradient dictionary less its approximation cost. If $q$ belongs to that
dictionary, the critic weakly reduces this risk, strictly whenever
$\operatorname{tr}\{(I-P)\Sigma\}>0$. The same fitted scalar
function and its labels may be reused for every current task having
the common continuation in Proposition~\ref{prop:r16menucache}.
\end{proposition}

\begin{proof}
Write $Pz-q=-(I-P)q+P\varepsilon$. Orthogonality removes the
cross term, and $P=P^\top=P^2$ gives
$\E\|P\varepsilon\|^2=\operatorname{tr}(P\Sigma)$.
The raw-label risk is $\operatorname{tr}(\Sigma)$, proving the
identity and its consequences. Neither step uses a Gaussian law or
independence of the coordinates of a label. Common continuation makes
the target gradient and its scalar feature representation independent
of the current task; evaluating an already fitted function for another
task therefore does not require fitting it again.
\end{proof}

The identity isolates a possible statistical contribution of learning
the continuation. It also specifies its price: a gradient dictionary
can impose enough approximation error to outweigh its variance
reduction. The twelve retained random-feature diagnostics and the
nonlinear comparisons address this tradeoff; the identity does not
determine their empirical signs. Its exact scope is a dictionary fixed
independently of the labels used for this least-squares fit and risk
at the declared arguments. The fully trained nonlinear NBO critic is
assessed by the independent economic comparisons below, rather than
being identified with that linear projection. A Raw procedure that
uses the same projection receives the same risk calculation. Task
reuse eliminates repeated fitting of a common object, while evaluation,
optimization, and any new state recursions remain part of each method's
recorded work.

\subsection{A direct conditional payoff statistic}
\label{app:r16menustatistic}

For one fixed query and two already selected actions $a,b$, let
$\overline F(x)=\{F(x;\xi)+F(x;-\xi)\}/2$, where a new base array
$\xi$ is independent of fitting and all other confirmation pairs. Put
$\gamma=2$ times the stored terminal-dispersion penalty and define
\[
 q_0^{ab}=-\frac\gamma2
       \{\|\Pi x_a\|_d^2-\|\Pi x_b\|_d^2\},\qquad
 R^{ab}=\overline F(x_a)-\overline F(x_b)-q_0^{ab}.
\]
The current reward and $q_0^{ab}$ are known conditional on this query
and its actions. Thus
\begin{equation}
 Q_\theta(y,a)-Q_\theta(y,b)
       =r_\theta(y,a)-r_\theta(y,b)+q_0^{ab}+\E R^{ab}.
 \label{eq:r16menucentered}
\end{equation}
The implementation may center at an enclosed binary64 representative of
$q_0^{ab}$. Its exact discrepancy is included in the range; the same
representative is added back to the known part of the payoff.

Let $N$ denote the number of dates, $n=N-1$, and let $\nu_I,\nu_C$
denote the stored idiosyncratic and common innovation amplitudes per
date. Set
\[
 H=(1+h\beta)^n,\qquad
 P=\beta_\psi\sum_{j=1}^{N-1}B_j(1+h\beta)^{j-1},\qquad
 D=\|x_a-x_b\|_d,
\]
and $s=\max\{\|\Pi x_a\|_d,\|\Pi x_b\|_d\}$. The following
coefficients contain no learned parameter:
\begin{align}
 C&=P+\gamma H\kappa nh
       +\gamma(H-1)\{s+\nu_I\sqrt{N(1-1/d)}\},\notag\\
 A&=\gamma(H-1)\nu_I\sqrt{N/d}.
 \label{eq:r16menuconstants}
\end{align}

\begin{proposition}[Query-conditional antithetic range]
\label{prop:r16menurange}
For one independent antithetic pair,
\[
 \Pr\{|R^{ab}|>D(C+Av)\}\leq e^{-v^2/2},\qquad v\geq0.
\]
Consequently, clipping at the prospectively fixed $v_*>0$ admits
the expectation allowance $DAe^{-v_*^2/2}/v_*$. These statements
hold conditional on the selected actions, so their geometry may enter
the range before the independent pair is generated.
\end{proposition}

\begin{proof}
Apply the pathwise mean-value identity along the straight segment from
$x_b$ to $x_a$. The reference adjoint has norm growth at most $H$;
its accumulated production component has norm at most $P$. Subtracting
the known terminal derivative $-\gamma\Pi x$ leaves the production
adjoint, a bounded accumulated-production term, and
$(J^\top-I)$ times the centered starting state and centered future
innovation. The unmultiplied terminal innovation cancels exactly in an
antithetic pair. The starting state is fixed. The norm of the projected
future innovation has mean at most $\nu_I\sqrt{N(1-1/d)}$ and,
as a function of its standard Gaussian array, normalized Lipschitz
constant at most $\nu_I\sqrt{N/d}$. Gaussian concentration gives
one event, uniform over the entire straight segment. Integration proves
the range, and integration of its tail proves the expectation allowance.
There is no prefix-state event because the query is conditioned upon.
\end{proof}

One observation is one antithetic pair, not its two dependent paths.
Different query-pair cells use independent arrays. Equal pair counts
at every query and at every declared training stream identify their
uniform finite average. For each contrast, the report pools only its
independent pair observations, clips each query at its proved range,
uses the maximum range in the independent nonidentical empirical-Bernstein
bound, and averages numerical and tail allowances. Dependence across
methods is retained by forming the direct contrast before inference.
The primary family comprises 216 two-sided events and uses probability
budget $0.018$. The separate scalar-accuracy family below uses $0.002$.
Their union remains within the fixed total $0.02$ budget.

\paragraph*{Arithmetic and Gaussian clipping.}
The numerical verifier uses the preserved protected polynomial for
$\tanh$, outward coefficient operations, and forward state-error
recursions with multiplier $1+h\beta$. Fixed dyadic spectral proposals
for the two declared coupling matrices are each certified by a fresh
outward LDL calculation. Postdecision representatives are interval
midpoints, so reconstructing their error accounts does not depend on
the last bits of a numerical singular-value or matrix-product routine.
A deterministic state cap
is computed from the fixed query, bounded drift and reference actions,
and the executed Gaussian cap. Production error is bounded by
$\beta_\psi$ times the state error plus its protected-function and
dot-product allowances. A terminal state error $e$ at a cap $M$
contributes at most $(\gamma/2)e(2M+e)$, before the separately charged
scalar reductions. The source records every component of this account.

The finite economy retains untruncated Gaussian innovations. For cap $z$,
\[
 \E\{G-\operatorname{clip}_{z}(G)\}^2
       \leq 4\phi(z)/z^3.
\]
This follows by writing the two-sided tail integral as
$2\phi(z)\int_0^\infty u^2e^{-zu-u^2/2}du$ and dropping
$e^{-u^2/2}$. Multiplication by $\nu_I^2+\nu_C^2$ and the
same state-error recursion gives an $L^2$ clipping allowance. The
production Lipschitz bound and conditional terminal Cauchy--Schwarz
then transfer the implemented expectation to the original finite
Gaussian economy. The clipping and arithmetic accounts are distinct.

\subsection{Accuracy over an entire scalar intervention class}
\label{app:r16menuaccuracy}

A separate comparison supplies an absolute reference on a restricted
decision class while retaining the high-dimensional state. At one fixed
state and current task, set $a(s)=a_L+s(a_R-a_L)$, $0\leq s\leq1$.
The endpoints are the common first-period reference consumption plus
or minus the fixed radius, intersected with the original task box.
All subsequent actions follow the same reference. Write
$q(s)=Q_\theta\{y,a(s)\}$ and $r=a_R-a_L$.

For the first and second variations of a future state with respect to
$s$, start with $V_1=h\|r\|_d$ and $W_1=0$. With the global
scalar and vector Hessian coefficients $H_m,H_x$ from the paired
transfer, recursively set
\[
 V_{j+1}=(1+h\beta)V_j,\qquad
 W_{j+1}=(1+h\beta)W_j+hH_xV_j^2.
\]
Let
\[
 S_N=\max\{\|\Pi x_{a_L}\|_d,\|\Pi x_{a_R}\|_d\}
          +\kappa(N-1)h+\nu_I\sqrt{N(1-1/d)}.
\]
Then an upper continuation curvature is
\[
 \Lambda=\sum_{j=1}^{N-1}B_j(H_mV_j^2+\beta_\psi W_j)
                                         +\gamma S_NW_N.
\]
For $u_i=\max(a_{L,i},a_{R,i})$, current utility weight $w$ and
adjustment coefficient $\eta$, put
\[
 \mu=A_0\left\{w\overline{r_i^2/u_i^2}+\eta\bar r^2\right\}-\Lambda.
\]

\begin{proposition}[Scalar Bellman curvature and a protected optimality gap]
\label{prop:r16menugap}
The preceding coefficients imply $q''(s)\leq-\mu$ throughout
$[0,1]$. If $[g_L,g_U]$ is a valid interval for $q'(s_0)$, then
\[
 \max_{0\leq s\leq1}q(s)-q(s_0)
 \leq\max_{\substack{g\in\{g_L,g_U\}\\0\leq s\leq1}}
           \{g(s-s_0)-\mu(s-s_0)^2/2\}.
\]
The maximum is a scalar quadratic calculation. It remains valid when
$\mu\leq0$; in that case it does not assert strong concavity.
\end{proposition}

\begin{proof}
Differentiating the finite transition recursion gives the displayed
variation bounds. The production chain rule contributes at most
$H_mV_j^2+\beta_\psi W_j$. The terminal second derivative is
$-\gamma\{\|\Pi v_N\|_d^2+\langle\Pi Y_N,\Pi w_N\rangle_d\}$.
Dropping its negative squared term and using the deterministic bound
on $w_N$ and the moment bound on $Y_N$ gives $\gamma S_NW_N$.
Subtracting the minimum curvature of current utility proves the
first claim. Taylor's integral remainder gives the quadratic support
at $s_0$. Maximizing that support over the gradient interval and
feasible scalar displacement proves the second.
\end{proof}

The gradient interval is obtained without identifying an automatic
derivative with an exact expectation. On a fixed neighboring interval
$[s_-,s_+]$ containing $s_0$, a fresh paired-payoff calculation bounds
$\{q(s_+)-q(s_-)\}/(s_+-s_-)$. A global absolute second-derivative
bound $L$ adds the deterministic allowance
\[
 \frac{L\{(s_0-s_-)^2+(s_+-s_0)^2\}}{2(s_+-s_-)}.
\]
This is the integral of the derivative's Lipschitz bound; it also
covers one-sided differences at a boundary. The declared scalar
candidate is selected from the fitted continuation before this bank.
Binary64 construction of the affine actions receives its own payoff
allowance. All 128 scalar events use 65,536 independent antithetic
pairs each and the fixed probability allocation. The resulting
reference covers every scalar action in the interval, including
unevaluated ones. Its scope is the specified finite decision class;
it is not a high-dimensional full-adapted-control or continuous-time
near-optimality assertion.
